\documentclass[12pt]{article}
\usepackage[utf8]{inputenc}
\usepackage[T1]{fontenc}
\usepackage[brazil]{babel}
\usepackage[a4paper,margin=2.5cm]{geometry}
\usepackage{enumitem}
\setlist[itemize]{noitemsep,topsep=4pt}
\title{Sprint 3 - S3\_04\\Análise de Algoritmos e Complexidade}
\author{Matheus Moreli Durão Silva\\Matrícula 12621CBS217\\Cibersegurança - Bacharelado - Noturno}
\date{2026}
\begin{document}
\maketitle

\section*{Entregável}
Relatório sobre Big-O e comparação de custos. Este documento registra o objetivo do card, a solução desenvolvida e a forma de conferência. Os arquivos de código e as evidências citadas permanecem na mesma pasta deste documento.

\section*{O que foi estudado e aplicado}
\begin{itemize}
\item O(1), O(log n), O(n), O(n log n) e O(n²)
\item Busca linear versus busca binária
\item Escalabilidade e escolha de solução
\item Justificativas com exemplos simples
\end{itemize}

\section*{Como conferir}
O relatório \texttt{comparacao\_complexidade.md} traz melhor, médio e pior casos de busca linear, busca binária, Bubble Sort e Insertion Sort, com justificativas e condições de uso. A busca binária exige dados ordenados; no inventário pequeno, a simplicidade da consulta por ID pode valer mais que ordenar os dados somente para buscar.

\section*{Conclusão}
O artefato foi organizado para demonstrar tanto o conceito estudado quanto sua aplicação prática. As evidências de execução só são consideradas válidas quando produzidas pela ferramenta correspondente.
\end{document}
